\documentclass[10pt]{article}
\usepackage[letterpaper,margin=0.62in]{geometry}
\usepackage{newtxtext,newtxmath}
\usepackage{microtype}
\usepackage{enumitem}
\usepackage{booktabs}
\usepackage[hidelinks]{hyperref}
\usepackage{titlesec}

\setlength{\parindent}{0pt}
\setlength{\parskip}{2.5pt}
\setlist[itemize]{leftmargin=1.1em,itemsep=0pt,topsep=1pt}
\setlist[enumerate]{leftmargin=1.35em,itemsep=0pt,topsep=1pt}
\titleformat{\section}{\bfseries\large}{}{0pt}{}[\vspace{-2pt}\titlerule]
\titlespacing*{\section}{0pt}{4pt}{2pt}

\title{\vspace{-1.35cm}\textbf{T-State Injection and Magic-State Resource Estimation}\\
\large QEC Simulation-(3) Pauli-Based Computation}
\author{Yongxiang Lyu \qquad Sheng Chen \qquad Zhuozi Xie}
\date{}

\begin{document}
\maketitle
\vspace{-1.1em}

\textbf{Course:} CSC/ECE 569 Quantum Computing \hfill
\textbf{Repository:} \href{https://github.com/cream-sketch/ECE-569}{github.com/cream-sketch/ECE-569}

\section{Motivation and Focus}

Stim and LightStim efficiently simulate stabilizer/QEC circuits containing Clifford operations, Pauli measurements, and Pauli noise, but Stim does not support non-Clifford gates such as T directly~\cite{gidney2021stim}. Pauli-based computation (PBC) motivates separating these two regimes: Clifford processing remains in the stabilizer/measurement domain, while the non-Clifford content appears as T-type magic-state resources~\cite{peres2022pbc}.

Rather than implementing a complete PBC compiler, this project focuses on two concrete problems:
\begin{enumerate}
\item \textbf{T-state preparation/injection:} replace an explicit T gate with a prepared magic state plus Clifford operations, measurement, and feed-forward.
\item \textbf{Resource estimation:} given a circuit T-count and target T-state error, estimate the number of distillation rounds and raw magic states required.
\end{enumerate}
This scope gives the project one implementation problem and one quantitative resource-analysis problem.

\section{Problem 1: T-State Preparation and T-Gate Injection}

The T gate is $T=\mathrm{diag}(1,e^{i\pi/4})$, and its magic state is
\[
|T\rangle=T|+\rangle=
\frac{|0\rangle+e^{i\pi/4}|1\rangle}{\sqrt{2}}.
\]
Preparing $|T\rangle$ by simply applying an ideal T gate is useful as a mathematical reference, but it does not solve the fault-tolerant problem because it assumes the non-Clifford operation that must ultimately be supplied. Fault-tolerant schemes instead prepare this resource offline and consume it through magic-state injection~\cite{bravyi2005magic}.

We will implement a T-injection gadget using a data qubit $|\psi\rangle$ and an ancilla $|T\rangle$:
\begin{enumerate}
\item prepare the ancilla in $|T\rangle$;
\item apply $\mathrm{CNOT}(\mathrm{data}\rightarrow\mathrm{ancilla})$;
\item measure the ancilla in the Z basis;
\item for measurement outcome 1, apply an S correction to the data qubit.
\end{enumerate}
For outcome 0, the data becomes $T|\psi\rangle$. For outcome 1, the uncorrected branch is $T^\dagger|\psi\rangle$ up to global phase; the Clifford S correction yields $T|\psi\rangle$. Thus, once the magic state is supplied, the online gadget uses only Clifford operations, measurement, and feed-forward.

\textbf{Implementation.} We will use Qiskit/statevector simulation to generate the ideal $|T\rangle$ reference and verify the injection gadget. Stim/LightStim will represent the surrounding Clifford operations, measurements, Pauli-frame corrections, and QEC circuitry. We will treat the magic state itself as the boundary to the non-Clifford model rather than claim that Stim directly represents its non-stabilizer amplitudes.

\textbf{Correctness experiment.} For $|0\rangle$, $|1\rangle$, $|+\rangle$, $|-\rangle$, and several random one-qubit inputs, we will compare a direct ideal T gate with injection. Output-state fidelity is the main metric; in the noiseless case the target is unit fidelity up to numerical precision. We will verify both measurement branches and test a small multi-T circuit to confirm that multiple injection gadgets compose correctly.

\section{Problem 2: Magic-State Distillation Resource Estimation}

After PBC or another Clifford+T analysis identifies the T resources, the next problem is their production cost. We will implement a resource estimator based on a single standard protocol, \textbf{15-to-1 T-state distillation}, rather than model every possible factory.

For sufficiently small input error $p$, a leading-order model is~\cite{bravyi2005magic}
\[
p_{\mathrm{out}}\approx 35p^3.
\]
For repeated rounds,
\[
p_{r+1}\approx 35p_r^3.
\]
Ignoring rejection/retry overhead in the first-order model, each round multiplies raw-state demand by 15:
\[
N_{\mathrm{raw/out}}\approx 15^r.
\]

The estimator inputs are circuit T-count $N_T$, raw-state error $p_{\mathrm{raw}}$, and target error $p_{\mathrm{target}}$. Optionally, an algorithm-level failure budget $\epsilon$ can define a simple per-T target such as $\epsilon/N_T$. Outputs will include the minimum number of rounds, error after each round, raw states per output, total raw states $N_T15^r$, and the number of high-fidelity T states consumed. With simple timing/parallelism assumptions, we will also report an estimated latency or throughput.

For example, $p_{\mathrm{raw}}=10^{-2}$ gives approximately $3.5\times10^{-5}$ after one idealized 15-to-1 round. A second round suppresses the error much further but raises the ideal raw-state cost from 15 to 225 states per final output. This directly exposes the fidelity-versus-resource tradeoff.

\section{Evaluation, Scope, and Deliverables}

\textbf{Experiment A -- Injection correctness.} We will report direct-T vs.\ injected-T fidelity, correctness of both measurement/feed-forward branches, Clifford gate count, measurement/correction count, and magic states consumed per T operation.

\textbf{Experiment B -- Resource scaling.} We will sweep representative T-counts $N_T=10,50,100,500,1000$, raw-state errors such as $10^{-2}$ and $10^{-3}$, and multiple target-error levels. Plots will show distillation rounds vs.\ target error, raw states vs.\ T-count, raw states vs.\ target fidelity, and total magic-state cost.

The project intentionally excludes a complete PBC compiler, full surface-code scheduling, lattice surgery, and a detailed physical magic-state factory. PBC provides the interface: it isolates T resources, and our work starts at that boundary.

Minimum deliverables are: (1) T-state/injection code and correctness tests; (2) a 15-to-1 resource estimator; (3) small Clifford+T integration examples; (4) resource-scaling plots/tables; and (5) documented limitations.

\section{Timeline}

\begin{center}
\small
\begin{tabular}{@{}ll@{}}
\toprule
\textbf{Period} & \textbf{Planned Work} \\
\midrule
Week 1 & Review Stim/LightStim, T injection, and 15-to-1 distillation \\
Week 2 & Implement ideal $|T\rangle$ reference and injection; verify both branches \\
Week 3 & Test multi-T examples and connect Clifford/measurement portions to Stim/LightStim \\
Week 4 & Implement 15-to-1 error/cost recurrence and resource estimator \\
Week 5 & Sweep T-count/error targets and generate resource plots \\
Week 6 & Integrate results, document limitations, and prepare final report \\
\bottomrule
\end{tabular}
\end{center}

{\small
\bibliographystyle{IEEEtran}
\bibliography{references}
}

\end{document}
